%============================================================
% CH1 · LESSON 4 — A FIRST TRAVERSAL
%============================================================
\subsection{A First Traversal}
\label{ssec:ch1-first-traversal}

The most basic question you can ask a graph is \textbf{reachability}: \emph{starting from vertex $s$, which vertices can I get to?} Two strategies answer it, and between them they are the backbone of half of this book:
\begin{itemize}[leftmargin=*]
    \item \textbf{Depth-first search (DFS)} --- walk as far as you can, and back up only when you are stuck. Like exploring a maze with one hand on the wall.
    \item \textbf{Breadth-first search (BFS)} --- explore in rings: first everything one step away, then everything two steps away, and so on. Like a rumour spreading, or ink spreading in water.
\end{itemize}
Both need one extra ingredient, a \texttt{visited} array, and both run in $\bigO{n+m}$. This lesson runs them by hand on the same small graph $G_3$; Chapter~2 (DFS) and Chapter~3 (BFS) study each one in depth.

% local drawing of G3: a pic whose vertex / edge styles are overridden per snapshot
\tikzset{
    g1l-s1/.style=gv, g1l-s2/.style=gv, g1l-s3/.style=gv, g1l-s4/.style=gv,
    g1l-s5/.style=gv, g1l-s6/.style=gv, g1l-s7/.style=gv,
    g1l-e12/.style=ge, g1l-e14/.style=ge, g1l-e23/.style=ge, g1l-e25/.style=ge, g1l-e45/.style=ge, g1l-e56/.style=ge,
    g1l-g3/.pic={
        \node[g1l-s1] (v1) at (0,0) {1};
        \node[g1l-s2] (v2) at (1.4,0.9) {2};
        \node[g1l-s4] (v4) at (1.4,-0.9) {4};
        \node[g1l-s3] (v3) at (2.9,0.9) {3};
        \node[g1l-s5] (v5) at (2.9,-0.9) {5};
        \node[g1l-s6] (v6) at (4.3,-0.9) {6};
        \node[g1l-s7] (v7) at (4.3,0.9) {7};
        \draw[g1l-e12] (v1) -- (v2); \draw[g1l-e14] (v1) -- (v4); \draw[g1l-e23] (v2) -- (v3);
        \draw[g1l-e25] (v2) -- (v5); \draw[g1l-e45] (v4) -- (v5); \draw[g1l-e56] (v5) -- (v6);
    }
}

\begin{diagrambox}
\begin{tikzpicture}[x=1.15cm, y=1cm]
    \pic {g1l-g3};
    \node[draw=darkgray, fill=lightgray!60, rounded corners=2pt, font=\ttfamily\small, align=left, inner sep=4pt, anchor=west] at (6.0,0)
        {7 6\\1 2\\1 4\\2 3\\2 5\\4 5\\5 6\\1};
    \node[font=\scriptsize\color{darkgray}, anchor=west] at (6.0,2.0) {input ($s$ last)};
    \node[note, anchor=west, text width=4.2cm] at (7.4,0) {adj[1] = 2 4\\adj[2] = 1 3 5\\adj[3] = 2\\adj[4] = 1 5\\adj[5] = 2 4 6\\adj[6] = 5\\adj[7] = (empty)};
\end{tikzpicture}
\figcap{$G_3$: $n=7$, $m=6$, start $s=1$. Vertex 7 is isolated, so the answer to ``what can 1 reach?'' is $\{1,2,3,4,5,6\}$. Neighbours are visited in the order of the adjacency lists (input order), which fixes the exact order of both searches.}
\end{diagrambox}

%------------------------------------------------------------
\subsubsection*{Why a visited array is essential}

The naive idea ``from $u$, go to every neighbour, and from there to every neighbour, \ldots'' never stops on a graph with a cycle: from 1 go to 2, from 2 back to 1, from 1 to 2 again, \ldots{} Even on a graph \emph{without} cycles it can explore the same vertex many times. The fix is one boolean per vertex: \textbf{mark a vertex the first time you reach it, and never enter a marked vertex again.}

\begin{diagrambox}
\begin{tikzpicture}[x=1cm, y=1cm]
    \begin{scope}
        \node[gv] (a) at (0,0.6) {1};
        \node[gv] (b) at (1.4,0.6) {2};
        \node[gv] (c) at (0.7,-0.6) {3};
        \draw[ge] (a) -- (b); \draw[ge] (b) -- (c); \draw[ge] (c) -- (a);
        \node[font=\scriptsize\bfseries] at (0.7,-1.5) {a triangle};
    \end{scope}
    \begin{scope}[xshift=3.4cm, yshift=0.9cm]
        \node[gvcur, minimum size=6mm] (r) at (1.8,0) {1};
        \node[gv, minimum size=6mm] (r2) at (0.9,-0.8) {2};
        \node[gv, minimum size=6mm] (r3) at (2.7,-0.8) {3};
        \node[gvB, minimum size=6mm] (r21) at (0.4,-1.6) {1};
        \node[gv, minimum size=6mm] (r23) at (1.4,-1.6) {3};
        \node[gvB, minimum size=6mm] (r31) at (2.2,-1.6) {1};
        \node[gv, minimum size=6mm] (r32) at (3.2,-1.6) {2};
        \foreach \a/\b in {r/r2, r/r3, r2/r21, r2/r23, r3/r31, r3/r32} \draw[faded] (\a) -- (\b);
        \node[font=\scriptsize\color{recursionred}] at (1.8,-2.25) {$\cdots$ forever};
        \node[font=\scriptsize\bfseries] at (1.8,-2.75) {without \texttt{visited}};
    \end{scope}
    \begin{scope}[xshift=8.6cm, yshift=0.9cm]
        \node[gvdone, minimum size=6mm] (r) at (0.9,0) {1};
        \node[gvdone, minimum size=6mm] (r2) at (0.9,-0.8) {2};
        \node[gvdone, minimum size=6mm] (r3) at (0.9,-1.6) {3};
        \draw[gtree] (r) -- (r2); \draw[gtree] (r2) -- (r3);
        \node[font=\scriptsize, anchor=west, text width=3cm] at (1.4,-0.8) {3 calls in total; every other edge leads to a marked vertex and is skipped};
        \node[font=\scriptsize\bfseries] at (1.8,-2.75) {with \texttt{visited}};
    \end{scope}
\end{tikzpicture}
\figcap{On a triangle the unmarked search revisits vertex 1 (red) again and again --- an infinite recursion. With marks, each vertex is entered once.}
\end{diagrambox}

%------------------------------------------------------------
\subsubsection*{Depth-first search (recursive)}

\begin{algorithmbox}[DFS: go deep first]
\texttt{dfs(u)}: mark $u$ as visited; for each neighbour $v$ of $u$: if $v$ is not visited, call \texttt{dfs(v)}.

When \texttt{dfs(v)} returns, the search is \emph{back at $u$} and tries $u$'s next neighbour. The chain of calls currently open (the recursion stack) is always a path from $s$ to the current vertex.
\end{algorithmbox}

\begin{lstlisting}[style=cpp]
#include <bits/stdc++.h>
using namespace std;

// Recursive DFS: which vertices can be reached from s?
// Input: n m, then m undirected edges "u v", then s.
vector<vector<int>> adj;
vector<bool> visited;
vector<int> order;                 // the order in which vertices are entered

void dfs(int u) {
    visited[u] = true;             // mark FIRST, then explore
    order.push_back(u);
    for (int v : adj[u])
        if (!visited[v])
            dfs(v);                // go deeper; come back when v is finished
}

int main() {
    int n, m;
    if (!(cin >> n >> m)) return 0;
    adj.assign(n + 1, vector<int>());
    for (int i = 0; i < m; i++) {
        int u, v;
        cin >> u >> v;
        adj[u].push_back(v);
        adj[v].push_back(u);
    }
    int s;
    cin >> s;
    visited.assign(n + 1, false);
    dfs(s);
    cout << "DFS order:";
    for (int v : order) cout << " " << v;
    cout << "\nunreachable:";
    for (int v = 1; v <= n; v++)
        if (!visited[v]) cout << " " << v;
    cout << "\n";
    return 0;
}
\end{lstlisting}

\begin{lstlisting}[style=trace]
input : the G3 input above
output: DFS order: 1 2 3 5 4 6
        unreachable: 7
\end{lstlisting}

\begin{tracebox}[DFS on $G_3$ from $s=1$, step by step]
Colours: \legendchip{gvcur}{current vertex}, \legendchip{gvseen}{on the stack (entered, not finished)}, \legendchip{gvdone}{finished}; thick blue edges are the \emph{tree edges} along which DFS entered a new vertex.

\begin{center}
\begin{tikzpicture}[x=0.62cm, y=0.62cm, every node/.append style={transform shape=false}]
    \tikzset{gv/.append style={minimum size=6mm, font=\scriptsize\bfseries}}
    % 1
    \begin{scope}[g1l-s1/.style=gvcur]
        \pic {g1l-g3};
        \node[font=\scriptsize, anchor=west] at (-0.3,-1.9) {\textbf{1.} enter 1};
        \node[font=\scriptsize\ttfamily, anchor=west] at (-0.3,-2.6) {stack: 1};
    \end{scope}
    % 2
    \begin{scope}[xshift=4.6cm, g1l-s1/.style=gvseen, g1l-s2/.style=gvcur, g1l-e12/.style=gtree]
        \pic {g1l-g3};
        \node[font=\scriptsize, anchor=west] at (-0.3,-1.9) {\textbf{2.} 1$\to$2 new: enter 2};
        \node[font=\scriptsize\ttfamily, anchor=west] at (-0.3,-2.6) {stack: 1 2};
    \end{scope}
    % 3
    \begin{scope}[xshift=9.2cm, g1l-s1/.style=gvseen, g1l-s2/.style=gvseen, g1l-s3/.style=gvcur,
                  g1l-e12/.style=gtree, g1l-e23/.style=gtree]
        \pic {g1l-g3};
        \node[font=\scriptsize, anchor=west] at (-0.3,-1.9) {\textbf{3.} 2$\to$1 seen; 2$\to$3: enter 3};
        \node[font=\scriptsize\ttfamily, anchor=west] at (-0.3,-2.6) {stack: 1 2 3};
    \end{scope}
    % 4
    \begin{scope}[yshift=-4.4cm, g1l-s1/.style=gvseen, g1l-s2/.style=gvseen, g1l-s3/.style=gvdone, g1l-s5/.style=gvcur,
                  g1l-e12/.style=gtree, g1l-e23/.style=gtree, g1l-e25/.style=gtree]
        \pic {g1l-g3};
        \node[font=\scriptsize, anchor=west] at (-0.3,-1.9) {\textbf{4.} 3 done, back to 2; enter 5};
        \node[font=\scriptsize\ttfamily, anchor=west] at (-0.3,-2.6) {stack: 1 2 5};
    \end{scope}
    % 5
    \begin{scope}[xshift=4.6cm, yshift=-4.4cm, g1l-s1/.style=gvseen, g1l-s2/.style=gvseen, g1l-s3/.style=gvdone, g1l-s5/.style=gvseen, g1l-s4/.style=gvcur,
                  g1l-e12/.style=gtree, g1l-e23/.style=gtree, g1l-e25/.style=gtree, g1l-e45/.style=gtree]
        \pic {g1l-g3};
        \node[font=\scriptsize, anchor=west] at (-0.3,-1.9) {\textbf{5.} 5$\to$2 seen; 5$\to$4: enter 4};
        \node[font=\scriptsize\ttfamily, anchor=west] at (-0.3,-2.6) {stack: 1 2 5 4};
    \end{scope}
    % 6
    \begin{scope}[xshift=9.2cm, yshift=-4.4cm, g1l-s1/.style=gvseen, g1l-s2/.style=gvseen, g1l-s3/.style=gvdone, g1l-s5/.style=gvseen, g1l-s4/.style=gvdone, g1l-s6/.style=gvcur,
                  g1l-e12/.style=gtree, g1l-e23/.style=gtree, g1l-e25/.style=gtree, g1l-e45/.style=gtree, g1l-e56/.style=gtree]
        \pic {g1l-g3};
        \node[font=\scriptsize, anchor=west] at (-0.3,-1.9) {\textbf{6.} 4: both seen, done; enter 6};
        \node[font=\scriptsize\ttfamily, anchor=west] at (-0.3,-2.6) {stack: 1 2 5 6};
    \end{scope}
\end{tikzpicture}
\end{center}

After step 6: vertex 6 sees only 5 (visited) and finishes; then 5, 2 finish; back at 1, the neighbour 4 is already visited, and 1 finishes. Vertex 7 is never touched. Entry order \textbf{1, 2, 3, 5, 4, 6}; maximum stack depth 4 (\texttt{1 2 5 4}).
\end{tracebox}

Notice what happened at step~5: DFS reached 4 through the long way round (1--2--5--4), not via the direct edge 1--4. DFS finds \emph{a} path to every reachable vertex, but not necessarily the shortest one.

%------------------------------------------------------------
\subsubsection*{Breadth-first search}

\begin{algorithmbox}[BFS: explore in rings]
Put $s$ in a queue with $\dist{s}=0$. Repeatedly take the vertex $u$ at the \emph{front} of the queue; for each neighbour $v$ that has not been discovered, set $\dist{v}=\dist{u}+1$ and push $v$ at the \emph{back}.

Because the queue is first-in first-out, all vertices at distance $d$ leave the queue before any vertex at distance $d+1$. So \texttt{dist[v]} is the \textbf{number of edges on a shortest path} from $s$ to $v$ --- a bonus DFS does not give.
\end{algorithmbox}

\begin{lstlisting}[style=cpp]
#include <bits/stdc++.h>
using namespace std;

// BFS: reachability from s plus the number of edges on a shortest path.
// Input: n m, then m undirected edges "u v", then s.
int main() {
    int n, m;
    if (!(cin >> n >> m)) return 0;
    vector<vector<int>> adj(n + 1);
    for (int i = 0; i < m; i++) {
        int u, v;
        cin >> u >> v;
        adj[u].push_back(v);
        adj[v].push_back(u);
    }
    int s;
    cin >> s;
    vector<int> dist(n + 1, -1);   // -1 = not discovered yet (doubles as visited)
    queue<int> q;
    dist[s] = 0;
    q.push(s);
    cout << "BFS order:";
    while (!q.empty()) {
        int u = q.front();
        q.pop();
        cout << " " << u;
        for (int v : adj[u]) {
            if (dist[v] == -1) {           // first time we see v
                dist[v] = dist[u] + 1;     // mark when PUSHED, not when popped
                q.push(v);
            }
        }
    }
    cout << "\ndist:";
    for (int v = 1; v <= n; v++) cout << " " << dist[v];
    cout << "\n";
    return 0;
}
\end{lstlisting}

\begin{lstlisting}[style=trace]
input : the G3 input above
output: BFS order: 1 2 4 3 5 6
        dist: 0 1 2 1 2 3 -1
\end{lstlisting}

\begin{tracebox}[BFS on $G_3$ from $s=1$, step by step]
Colours: \legendchip{gvcur}{popped now}, \legendchip{gvseen}{in the queue}, \legendchip{gvdone}{popped earlier}; green numbers are \texttt{dist}; thick blue edges are the edges along which a vertex was discovered.

\begin{center}
\begin{tikzpicture}[x=0.62cm, y=0.62cm]
    \tikzset{gv/.append style={minimum size=6mm, font=\scriptsize\bfseries}}
    % 1
    \begin{scope}[g1l-s1/.style=gvcur, g1l-s2/.style=gvseen, g1l-s4/.style=gvseen, g1l-e12/.style=gtree, g1l-e14/.style=gtree]
        \pic {g1l-g3};
        \foreach \v/\d in {1/0, 2/1, 4/1} \node[gdist, above left=-1pt of v\v] {\d};
        \node[font=\scriptsize, anchor=west] at (-0.3,-1.9) {\textbf{1.} pop 1: push 2, 4};
        \node[font=\scriptsize\ttfamily, anchor=west] at (-0.3,-2.6) {queue: 2 4};
    \end{scope}
    % 2
    \begin{scope}[xshift=4.6cm, g1l-s1/.style=gvdone, g1l-s2/.style=gvcur, g1l-s4/.style=gvseen, g1l-s3/.style=gvseen, g1l-s5/.style=gvseen,
                  g1l-e12/.style=gtree, g1l-e14/.style=gtree, g1l-e23/.style=gtree, g1l-e25/.style=gtree]
        \pic {g1l-g3};
        \foreach \v/\d in {1/0, 2/1, 4/1, 3/2, 5/2} \node[gdist, above left=-1pt of v\v] {\d};
        \node[font=\scriptsize, anchor=west] at (-0.3,-1.9) {\textbf{2.} pop 2: 1 seen; push 3, 5};
        \node[font=\scriptsize\ttfamily, anchor=west] at (-0.3,-2.6) {queue: 4 3 5};
    \end{scope}
    % 3
    \begin{scope}[xshift=9.2cm, g1l-s1/.style=gvdone, g1l-s2/.style=gvdone, g1l-s4/.style=gvcur, g1l-s3/.style=gvseen, g1l-s5/.style=gvseen,
                  g1l-e12/.style=gtree, g1l-e14/.style=gtree, g1l-e23/.style=gtree, g1l-e25/.style=gtree]
        \pic {g1l-g3};
        \foreach \v/\d in {1/0, 2/1, 4/1, 3/2, 5/2} \node[gdist, above left=-1pt of v\v] {\d};
        \node[font=\scriptsize, anchor=west] at (-0.3,-1.9) {\textbf{3.} pop 4: 1, 5 seen};
        \node[font=\scriptsize\ttfamily, anchor=west] at (-0.3,-2.6) {queue: 3 5};
    \end{scope}
    % 4
    \begin{scope}[yshift=-4.4cm, g1l-s1/.style=gvdone, g1l-s2/.style=gvdone, g1l-s4/.style=gvdone, g1l-s3/.style=gvcur, g1l-s5/.style=gvseen,
                  g1l-e12/.style=gtree, g1l-e14/.style=gtree, g1l-e23/.style=gtree, g1l-e25/.style=gtree]
        \pic {g1l-g3};
        \foreach \v/\d in {1/0, 2/1, 4/1, 3/2, 5/2} \node[gdist, above left=-1pt of v\v] {\d};
        \node[font=\scriptsize, anchor=west] at (-0.3,-1.9) {\textbf{4.} pop 3: 2 seen};
        \node[font=\scriptsize\ttfamily, anchor=west] at (-0.3,-2.6) {queue: 5};
    \end{scope}
    % 5
    \begin{scope}[xshift=4.6cm, yshift=-4.4cm, g1l-s1/.style=gvdone, g1l-s2/.style=gvdone, g1l-s4/.style=gvdone, g1l-s3/.style=gvdone, g1l-s5/.style=gvcur, g1l-s6/.style=gvseen,
                  g1l-e12/.style=gtree, g1l-e14/.style=gtree, g1l-e23/.style=gtree, g1l-e25/.style=gtree, g1l-e56/.style=gtree]
        \pic {g1l-g3};
        \foreach \v/\d in {1/0, 2/1, 4/1, 3/2, 5/2, 6/3} \node[gdist, above left=-1pt of v\v] {\d};
        \node[font=\scriptsize, anchor=west] at (-0.3,-1.9) {\textbf{5.} pop 5: 2, 4 seen; push 6};
        \node[font=\scriptsize\ttfamily, anchor=west] at (-0.3,-2.6) {queue: 6};
    \end{scope}
    % 6
    \begin{scope}[xshift=9.2cm, yshift=-4.4cm, g1l-s1/.style=gvdone, g1l-s2/.style=gvdone, g1l-s4/.style=gvdone, g1l-s3/.style=gvdone, g1l-s5/.style=gvdone, g1l-s6/.style=gvcur, g1l-s7/.style=gvdead,
                  g1l-e12/.style=gtree, g1l-e14/.style=gtree, g1l-e23/.style=gtree, g1l-e25/.style=gtree, g1l-e56/.style=gtree]
        \pic {g1l-g3};
        \foreach \v/\d in {1/0, 2/1, 4/1, 3/2, 5/2, 6/3} \node[gdist, above left=-1pt of v\v] {\d};
        \node[font=\scriptsize, anchor=west] at (-0.3,-1.9) {\textbf{6.} pop 6: 5 seen};
        \node[font=\scriptsize\ttfamily, anchor=west] at (-0.3,-2.6) {queue: (empty) --- stop};
    \end{scope}
\end{tikzpicture}
\end{center}

Pop order \textbf{1, 2, 4, 3, 5, 6}; final \texttt{dist} $=(0,1,2,1,2,3,-1)$ for vertices $1..7$. Vertex 7 keeps $-1$: unreachable.
\end{tracebox}

\begin{diagrambox}
\begin{tikzpicture}[x=1.3cm, y=1cm]
    \foreach \d/\x in {0/0, 1/1.6, 2/3.2, 3/4.8} {
        \fill[primaryblue!\the\numexpr 6+5*\d\relax, rounded corners=4pt] (\x-0.55,-1.5) rectangle (\x+0.55,1.5);
        \node[font=\scriptsize\bfseries\color{deepblue}] at (\x,1.8) {dist $\d$};
    }
    \node[gvsrc] (1) at (0,0) {1};
    \node[gv] (2) at (1.6,0.7) {2};
    \node[gv] (4) at (1.6,-0.7) {4};
    \node[gv] (3) at (3.2,0.7) {3};
    \node[gv] (5) at (3.2,-0.7) {5};
    \node[gv] (6) at (4.8,-0.7) {6};
    \draw[gtree] (1) -- (2); \draw[gtree] (1) -- (4); \draw[gtree] (2) -- (3); \draw[gtree] (2) -- (5); \draw[gtree] (5) -- (6);
    \draw[gfaded, dashed] (4) -- (5);
    \node[gvdead] (7) at (6.4,0) {7};
    \node[font=\scriptsize\color{darkgray}] at (6.4,-0.7) {never reached};
\end{tikzpicture}
\figcap{BFS layers of $G_3$. Every edge joins two vertices whose distances differ by at most 1; the dashed edge 4--5 joins layer 1 to layer 2 but was not used to discover anyone.}
\end{diagrambox}

%------------------------------------------------------------
\subsubsection*{Same graph, two different trees}

Each search discovers every reachable vertex through exactly one edge --- its \textbf{tree edge}. The tree edges form a tree (the \textbf{DFS tree} or \textbf{BFS tree}) rooted at $s$ that spans everything reachable. The two trees differ:

\begin{diagrambox}
\begin{tikzpicture}[x=1cm, y=0.9cm]
    \begin{scope}
        \node[gvsrc] (1) at (0,0) {1};
        \node[gv] (2) at (0,-1) {2};
        \node[gv] (3) at (-0.9,-2) {3};
        \node[gv] (5) at (0.9,-2) {5};
        \node[gv] (4) at (0.3,-3) {4};
        \node[gv] (6) at (1.5,-3) {6};
        \draw[gtree] (1) -- (2); \draw[gtree] (2) -- (3); \draw[gtree] (2) -- (5); \draw[gtree] (5) -- (4); \draw[gtree] (5) -- (6);
        \draw[gback, -] (4) to[bend left=50] (1);
        \node[font=\small\bfseries] at (0.3,-3.9) {DFS tree: depth 3};
        \node[font=\scriptsize\color{darkgray}] at (0.3,-4.35) {non-tree edge 1--4 (dashed)};
    \end{scope}
    \begin{scope}[xshift=6cm]
        \node[gvsrc] (1) at (0,0) {1};
        \node[gv] (2) at (-0.8,-1) {2};
        \node[gv] (4) at (0.8,-1) {4};
        \node[gv] (3) at (-1.4,-2) {3};
        \node[gv] (5) at (-0.2,-2) {5};
        \node[gv] (6) at (-0.2,-3) {6};
        \draw[gtree] (1) -- (2); \draw[gtree] (1) -- (4); \draw[gtree] (2) -- (3); \draw[gtree] (2) -- (5); \draw[gtree] (5) -- (6);
        \draw[gback, -] (4) -- (5);
        \node[font=\small\bfseries] at (0,-3.9) {BFS tree: depth = distance};
        \node[font=\scriptsize\color{darkgray}] at (0,-4.35) {non-tree edge 4--5 (dashed)};
    \end{scope}
\end{tikzpicture}
\figcap{Both trees contain the same six vertices (everything reachable from 1) and $6-1=5$ edges. In the BFS tree the depth of a vertex equals its distance from $s$; in the DFS tree vertex 4 sits at depth 3 although $\dist{4}=1$.}
\end{diagrambox}

\begin{comparisonbox}[DFS vs.\ BFS at a glance]
\small
\begin{tabularx}{\linewidth}{@{}>{\raggedright\arraybackslash}p{4.2cm} X X@{}}
\toprule
 & \textbf{DFS} & \textbf{BFS} \\
\midrule
data structure & recursion (a stack) & a queue \\
order on $G_3$ & 1, 2, 3, 5, 4, 6 & 1, 2, 4, 3, 5, 6 \\
reachability & yes & yes \\
shortest paths (unweighted) & no & \textbf{yes}: \texttt{dist[]} \\
code length & shortest & a few lines more \\
deep graphs ($10^6$ vertices on a path) & recursion may overflow & safe \\
strengths later in the book & cycles, bridges, SCC, trees (Ch.~2) & distances, grids, states (Ch.~3) \\
\bottomrule
\end{tabularx}
\end{comparisonbox}

%------------------------------------------------------------
\subsubsection*{Why both run in $\bigO{n+m}$}

\begin{invariantbox}[Each vertex is expanded once, each edge is looked at twice]
A vertex is marked the moment it is first reached and is never entered (DFS) or pushed (BFS) again. So each vertex $u$ is \emph{expanded} --- its adjacency list scanned --- at most once, at a cost of $\deg(u)+\bigO{1}$. Summing over all vertices, the total work is
\[
    \sum_{u}\big(\deg(u)+1\big) \;=\; 2m + n \;=\; \bigO{n+m}
\]
by the handshake lemma (for a directed graph the sum of out-degrees is $m$). On $G_3$ both searches perform exactly $2m=12$ neighbour checks.
\end{invariantbox}

\begin{warningbox}[In BFS, mark when you push, not when you pop]
If you set \texttt{dist[v]} only when $v$ is popped, the same vertex can be pushed by several neighbours before it is first popped. The answers stay correct, but the queue can grow to $\bigO{m}$ entries and on dense graphs the running time explodes. Mark at push time, as in the code above.
\end{warningbox}

%------------------------------------------------------------
\subsubsection*{The recursion-depth trap}

The recursive DFS keeps one stack frame for every vertex on the current path. On a path graph 1--2--3--$\cdots$--$n$ with $n=10^6$ the recursion goes $10^6$ levels deep. Codeforces gives a large stack (256 MB) and this is usually fine there, but on a default 1--8 MB stack (Windows at home, many other judges, LeetCode) it is a \stackoverflow{} --- reported as a mysterious runtime error.

\begin{diagrambox}
\begin{tikzpicture}[x=1cm, y=1cm]
    \foreach \i/\l in {0/1, 1/2, 2/3, 3/4} \node[gvseen] (p\i) at (\i*1.3,0) {\l};
    \node (dots) at (5.2,0) {$\cdots$};
    \node[gvcur, minimum size=9mm, font=\scriptsize\bfseries] (pn) at (6.6,0) {$10^6$};
    \foreach \a/\b in {p0/p1, p1/p2, p2/p3, p3/dots, dots/pn} \draw[gtree] (\a) -- (\b);
    \foreach \i/\h in {0/0.35, 1/0.7, 2/1.05, 3/1.4}
        \fill[recursionred!25, draw=recursionred] (8.2,-0.6+\h-0.35) rectangle (10.2,-0.6+\h);
    \node[font=\scriptsize] at (9.2,1.15) {$\vdots$};
    \fill[recursionred!25, draw=recursionred] (8.2,1.35) rectangle (10.2,1.7);
    \node[font=\scriptsize\ttfamily] at (9.2,-0.42) {dfs(1)};
    \node[font=\scriptsize\ttfamily] at (9.2,-0.07) {dfs(2)};
    \node[font=\scriptsize\ttfamily] at (9.2,0.28) {dfs(3)};
    \node[font=\scriptsize\ttfamily] at (9.2,0.63) {dfs(4)};
    \node[font=\scriptsize\ttfamily] at (9.2,1.52) {dfs($10^6$)};
    \node[font=\scriptsize\bfseries\color{recursionred}] at (9.2,2.0) {$10^6$ open frames};
\end{tikzpicture}
\figcap{DFS on a long path: every call is still open when the next one starts.}
\end{diagrambox}

When the depth can exceed about $10^5$, either use BFS (if the order does not matter, as for plain reachability) or replace the recursion by an explicit stack. The iterative version below remembers, for each vertex on the stack, which neighbour to try next, so it visits vertices in \emph{exactly} the same order as the recursive one.

\begin{lstlisting}[style=cpp]
#include <bits/stdc++.h>
using namespace std;

// Iterative DFS with an explicit stack: no recursion, so no stack overflow
// even on a path of 10^6 vertices. It visits exactly the same vertices.
int main() {
    int n, m;
    if (!(cin >> n >> m)) return 0;
    vector<vector<int>> adj(n + 1);
    for (int i = 0; i < m; i++) {
        int u, v;
        cin >> u >> v;
        adj[u].push_back(v);
        adj[v].push_back(u);
    }
    int s;
    cin >> s;
    vector<bool> visited(n + 1, false);
    vector<size_t> it(n + 1, 0);     // it[u] = next neighbour index of u to try
    vector<int> st;                  // the "call stack", made by hand
    visited[s] = true;
    st.push_back(s);
    cout << "DFS order: " << s;
    while (!st.empty()) {
        int u = st.back();
        if (it[u] == adj[u].size()) { st.pop_back(); continue; }   // u finished
        int v = adj[u][it[u]++];
        if (!visited[v]) {
            visited[v] = true;
            cout << " " << v;
            st.push_back(v);                                       // "call" dfs(v)
        }
    }
    cout << "\n";
    return 0;
}
\end{lstlisting}

\begin{lstlisting}[style=trace]
input : the G3 input above
output: DFS order: 1 2 3 5 4 6          (identical to the recursive version)
\end{lstlisting}

Chapters~2 and~3 go much deeper: entry/exit times, edge classification, cycle detection and bridges for DFS; multi-source BFS, 0-1 BFS and BFS on states for BFS. For the rest of this chapter, the two programs above are all you need --- Patterns~2--4 are applications of them.

%------------------------------------------------------------
\subsubsection*{Check your understanding}

\begin{enumerate}[leftmargin=*]
    \item In $G_3$, change the input order so that the edge \texttt{1 4} comes before \texttt{1 2}. What is the new DFS order from 1?
    \item On a star (vertex 1 joined to $2,\dots,n$), what is the maximum recursion depth of DFS from 1? From 2?
    \item A graph has $n=10^5$ vertices and $m=0$ edges. How much work does a BFS from vertex 1 do? How much work does \emph{looping over all vertices and starting a BFS from every unvisited one} do?
    \item Why can BFS compute shortest paths, but DFS cannot? Answer with $G_3$.
    \item You run BFS and forget to set \texttt{dist[s] = 0} before the loop (it stays $-1$). What goes wrong?
\end{enumerate}

\begin{solutionbox}[Answers]
\begin{enumerate}[leftmargin=*]
    \item Then \texttt{adj[1] = 4 2}, and DFS goes $1\to4\to5\to2\to3$, back to 5, then $5\to6$: order \textbf{1, 4, 5, 2, 3, 6}. (Running the program on the reordered input prints exactly this.) A different but equally valid DFS.
    \item From 1: depth 1 (every leaf is a child of 1). From 2: depth 2 ($2\to1\to$ each other leaf).
    \item One BFS from 1 does $\bigO{1}$ work (no edges). Starting a search from every unvisited vertex costs $\bigO{n+m}$ in total, because each vertex is still expanded only once over all the searches --- this is how components are counted (Pattern~2, Chapter~4).
    \item BFS pops vertices in order of distance, so a vertex is first discovered along a shortest path. DFS discovered 4 along 1--2--5--4 (3 edges) although the edge 1--4 gives distance 1.
    \item Everything shifts and $s$ is processed twice. The neighbours of $s$ get $-1+1=0$; $s$ still looks undiscovered, so the first neighbour expanded pushes it again with distance 1. On $G_3$ the modified program prints \texttt{BFS order: 1 2 4 1 3 5 6} and \texttt{dist: 1 0 1 0 1 2 -1}. Always initialise the source before the loop.
\end{enumerate}
\end{solutionbox}
